% Generated draft: requires mathematical and source review.
\documentclass[11pt]{article}
\usepackage[margin=25mm]{geometry}
\usepackage{fontspec}
\setmainfont{Latin Modern Roman}
\usepackage{amsmath,amssymb,mathtools,unicode-math}
\setmathfont{Latin Modern Math}
\usepackage{enumitem,needspace,xurl,hyperref}
\setlength{\parindent}{0pt}
\setlength{\parskip}{5pt}
\setlength{\emergencystretch}{3em}
\newcommand{\N}{\mathbb{N}}
\newcommand{\Z}{\mathbb{Z}}
\newcommand{\Q}{\mathbb{Q}}
\newcommand{\R}{\mathbb{R}}
\newcommand{\C}{\mathbb{C}}
\newcommand{\F}{\mathbb{F}}
\newcommand{\A}{\mathbb{A}}
\newcommand{\PP}{\mathbb{P}}
\newcommand{\E}{\mathbb{E}}
\newcommand{\eps}{\varepsilon}
\newcommand{\dd}{\,\mathrm{d}}
\begin{document}
\section*{UnsolvedMath ID: 30005187}
\textbf{Source problem code:} catalogue entry
\subsection{Definitions and mathematical statement}
The supplied record contains a research question and a disrupted formula for a fluctuation, rather than a complete mathematical assertion. The readable formula can be given the following conditional interpretation. Let \(\mathcal A\) be an algebra represented on a Hilbert space \(\mathcal H\), with its represented elements denoted by the same symbols. Let \(D\) be a densely defined operator. For a finite positive integer \(n\), choose \(a_j,b_j\in\mathcal A\), \(1\leq j\leq n\), such that the commutators \([D,b_j]=Db_j-b_jD\) are defined on a common dense domain and extend to bounded operators, and such that
\[
V=\sum_{j=1}^{n}a_j[D,b_j]\in\mathcal B(\mathcal H)_{\mathrm{sa}}.
\]
Here \(\mathcal B(\mathcal H)_{\mathrm{sa}}\) denotes the bounded self-adjoint operators on \(\mathcal H\). These conditions explain how to read the displayed fragment; they do not recover the complete hypotheses of the contribution. The fragment also contains the notation \(\pi_D(A)\), but the object \(A\) and the map \(\pi_D\) are not defined in the supplied record.

The recoverable research question is: can the full spectral action be described field-theoretically, without restricting to its scale-invariant part? The record does not specify the spectral action functional, its admissible inputs, the meaning of a field-theoretic description, or a criterion for such a description to be valid. Thus no precise universally quantified conjecture or existence assertion can be recovered from this evidence.

There is also a substantive qualification: the supplied statement-verification field says that the contribution poses this question rhetorically and then answers it through its main spectral-action expansion theorem. It marks the statement as unrecoverable and promotes no separate clean statement. The supplied open label and broader literature assessment therefore do not establish that this particular question is an unresolved mathematical problem.
\subsection{Short English statement}
The record asks for a field-theoretic description of the entire spectral action, beyond its scale-invariant part, and includes a fragment describing a bounded self-adjoint fluctuation \(V\). Missing definitions and hypotheses prevent a precise conjecture. Moreover, the supplied verification says the contribution itself answers the rhetorical question, despite the registry's open label. Human review of the cited contribution is needed to identify any intended unresolved target.
\subsection{Sources}
\begin{itemize}
\item Supplied local UnsolvedMath registry record, permanent ID \(30005187\), title: Field-Theoretic Description of the Full Spectral Action; exact source problem code: \texttt{catalogue entry}. The record belongs to the Oberwolfach Reports Open Problems collection. No individual proposer is supplied; the proposed year is recorded as 2022.
\item Supplied source citation: Non-Commutative Geometry and Cyclic Homology (2023), Oberwolfach Reports, DOI 10.4171/owr/2022/35, pp. 1987--2058. Source URL and literature-source link labelled original Oberwolfach report: \url{https://doi.org/10.4171/owr/2022/35}.
\item Supplied literature-source link: Spectral Action in Noncommutative Geometry, \url{https://arxiv.org/abs/1902.05306}.
\item Literature checked at: \texttt{2026-08-22}. This is supplied dated curation, not an independently verified assessment. It labels the topic open and states that the spectral action principle and several asymptotic or scale-invariant sectors are established, while a general field-theoretic description of the full spectral action remains an open research direction. The background identifies this as literature triage derived from the supplied Oberwolfach curation and calls for independent verification of time-sensitive status and cited claims.
\item The supplied statement-verification field says the contribution answers the rhetorical question with its main spectral-action expansion theorem, promotes no separate clean statement, and retains the original record publicly for provenance. Its statement status is \texttt{unrecoverable}; its separate problem-status label is \texttt{open}. These are supplied registry qualifications, not verified mathematical conclusions.
\end{itemize}
\end{document}
